%
% teil2.tex -- Lösung mit verallgemeinerten Potenzreihen
%
% (c) 2026 Gian Cavegn, OST Ostschweizer Fachhochschule
%
% !TEX root = ../../buch.tex
% !TEX encoding = UTF-8
% !TeX spellcheck = de_CH
%
\section{Lösung mit verallgemeinerten Potenzreihen
\label{bessel:section:loesung}}
\kopfrechts{Lösung}
%
% Ziel dieses Abschnitts.
% Die allgemeine Theorie aus dem Fuchs-Abschnitt nicht wiederholen,
% sondern anwenden.
% Der Mehrwert liegt darin, dass die Bessel-Gleichung ein besonders
% sauberer Spezialfall ist.
% Es gilt F(varrho) = varrho^2 - nu^2, also varrho = +- nu.
%
% Notation bewusst wie im Buch, also varrho für den Exponenten der
% verallgemeinerten Potenzreihe, nu für die Ordnung.
%
Die Bessel-Gleichung
\begin{equation}
x^2y'' + xy' + (x^2-\nu^2)y = 0
\label{bessel:loesung:eqn:bessel}
\end{equation}
ist eine fuchssche Differentialgleichung.
Das in Abschnitt~\ref{buch:differentialgleichungen:section:fuchs} entwickelte Verfahren lässt sich deshalb direkt anwenden.
Wir wiederholen die allgemeine Theorie hier nicht, sondern setzen ein und beobachten, was dabei herauskommt.
Bemerkenswert ist, wie einfach der Spezialfall wird.
Alle Grössen der allgemeinen Rekursion nehmen für die Bessel-Gleichung eine besonders übersichtliche Form an.
 
%
% Die Bessel-Gleichung als fuchssche Gleichung
%
\subsection{Die Bessel-Gleichung als fuchssche Gleichung
\label{bessel:loesung:subsection:fuchssch}}
Der Vergleich von~\eqref{bessel:loesung:eqn:bessel} mit der allgemeinen Form \eqref{buch:differentialgleichungen:fuchs:eqn:dgl} liefert
\begin{equation}
p(x) = 1,
\qquad
q(x) = x^2-\nu^2.
\label{bessel:loesung:eqn:pq}
\end{equation}
Beide Funktionen sind Polynome, ihre Potenzreihenentwicklungen um $0$ brechen also nach wenigen Gliedern ab.
Für die Koeffizienten $p_l$ und $q_l$ bedeutet das
\begin{equation}
p_0 = 1,
\quad
p_l = 0
\ \text{für}\ l \ge 1,
\qquad
q_0 = -\nu^2,
\quad
q_1 = 0,
\quad
q_2 = 1,
\quad
q_l = 0
\ \text{für}\ l\ge 3.
\label{bessel:loesung:eqn:koeffizienten}
\end{equation}
Der Ansatz ist eine verallgemeinerte Potenzreihe im Sinne der Definition~\ref{buch:differentialgleichungen:verallgemeinert:def}, also
\begin{equation}
y(x)
=
\sum_{k=0}^\infty
a_k x^{k+\varrho}
=
x^\varrho
\sum_{k=0}^\infty
a_k x^k,
\qquad
a_0 \ne 0.
\label{bessel:loesung:eqn:ansatz}
\end{equation}
Der Faktor $x^\varrho$ ist der ganze Unterschied zu einem gewöhnlichen Potenzreihenansatz.
Er ist genau das, was das Verfahren rettet, denn er erlaubt der Lösung, sich bei $x=0$ wie eine beliebige, auch nicht ganzzahlige Potenz zu verhalten.
 
%
% Indexgleichung
%
\subsection{Die Indexgleichung
\label{bessel:loesung:subsection:index}}
Die Indexgleichung des Fuchs-Verfahrens lautet $F(\varrho) = \varrho(\varrho-1) + \varrho p_0 + q_0 = 0$.
Mit den Werten~\eqref{bessel:loesung:eqn:koeffizienten} wird daraus
\begin{equation}
F(\varrho)
=
\varrho(\varrho-1) + \varrho - \nu^2
=
\varrho^2 - \nu^2.
\label{bessel:loesung:eqn:index}
\end{equation}
Die beiden linearen Terme heben sich weg, und es bleibt die denkbar einfachste quadratische Gleichung.
Ihre Lösungen sind
\begin{equation}
\varrho_1 = \nu,
\qquad
\varrho_2 = -\nu.
\label{bessel:loesung:eqn:varrho}
\end{equation}
Der Exponent der verallgemeinerten Potenzreihe ist also nichts anderes als die Ordnung der Gleichung, mit beiden Vorzeichen.
Das ist keine blosse Rechenkuriosität, sondern lässt sich direkt ablesen.
Für kleine $x$ ist der Term $x^2$ in~\eqref{bessel:loesung:eqn:bessel}
gegenüber $\nu^2$ vernachlässigbar, und die Gleichung geht in
\[
x^2y'' + xy' - \nu^2 y = 0
\]
über.
Das ist eine eulersche Differentialgleichung, deren Lösungen genau $x^{\nu}$ und $x^{-\nu}$ sind.
Die Indexgleichung beschreibt also das Verhalten der Lösung in unmittelbarer Nähe der Singularität.
 
%
% Rekursionsformel
%
\subsection{Die Rekursionsformel für die Koeffizienten
\label{bessel:loesung:subsection:rekursion}}
Setzt man den Ansatz~\eqref{bessel:loesung:eqn:ansatz} in~\eqref{bessel:loesung:eqn:bessel} ein, so liefert der Koeffizientenvergleich die Bedingungen \eqref{buch:differentialgleichungen:fuchs:eqn:bedingungen}.
Für die Bessel-Gleichung vereinfachen sie sich drastisch, weil von den $p_l$ und $q_l$ nur $p_0$, $q_0$ und $q_2$ von null verschieden sind.
Der Koeffizient von $x^{n+\varrho}$ ist
\begin{equation}
\bigl((n+\varrho)^2-\nu^2\bigr)a_n + a_{n-2} = 0,
\qquad n \ge 2,
\label{bessel:loesung:eqn:rekursion}
\end{equation}
wobei der Term $a_{n-2}$ genau von $q_2=1$ herrührt.
% Kontrolle. F(n+varrho) = (n+varrho)(n+varrho-1) + (n+varrho) - nu^2
%            = (n+varrho)^2 - nu^2. Stimmt mit (bessel:loesung:eqn:index).
Für $n=1$ bleibt
\begin{equation}
\bigl((1+\varrho)^2-\nu^2\bigr)a_1 = 0.
\label{bessel:loesung:eqn:a1}
\end{equation}
 
Wir wählen zunächst $\varrho=\nu$ mit $\nu\ge 0$.
Dann ist $(1+\nu)^2-\nu^2 = 1+2\nu > 0$, aus \eqref{bessel:loesung:eqn:a1} folgt also $a_1=0$.
Wegen~\eqref{bessel:loesung:eqn:rekursion} verschwinden damit sämtliche Koeffizienten mit ungeradem Index,
\begin{equation}
a_1 = a_3 = a_5 = \dots = 0.
\label{bessel:loesung:eqn:ungerade}
\end{equation}
Die Reihe enthält nur gerade Potenzen.
 
Für die geraden Koeffizienten liefert \eqref{bessel:loesung:eqn:rekursion} mit $\varrho=\nu$
\begin{equation}
a_n
=
-\frac{a_{n-2}}{(n+\nu)^2-\nu^2}
=
-\frac{a_{n-2}}{n(n+2\nu)}.
\label{bessel:loesung:eqn:rekursionnu}
\end{equation}
Schreibt man $n=2k$, so wird daraus
\[
a_{2k}
=
-\frac{a_{2k-2}}{2k\,(2k+2\nu)}
=
-\frac{a_{2k-2}}{4k(k+\nu)},
\]
und durch wiederholtes Einsetzen ergibt sich die geschlossene Form
\begin{equation}
a_{2k}
=
\frac{(-1)^k a_0}{4^k\, k!\,(\nu+1)(\nu+2)\cdots(\nu+k)}.
\label{bessel:loesung:eqn:a2k}
\end{equation}
Der Nenner wächst wie $4^k(k!)^2$, die Reihe konvergiert daher für alle $x\in\mathbb{R}$.

% [TODO: Konvergenz sauber mit dem Quotientenkriterium ausführen,
%        Verweis auf den Konvergenzradius-Abschnitt im Buch]
 
%
% Die Bessel-Funktion erster Art
%
\subsection{Die Bessel-Funktion erster Art
\label{bessel:loesung:subsection:jnu}}
Der Koeffizient $a_0$ ist frei wählbar.
Die übliche Normierung ist
\begin{equation}
a_0
=
\frac{1}{2^\nu\,\Gamma(\nu+1)},
\label{bessel:loesung:eqn:normierung}
\end{equation}
und sie ist nicht willkürlich.
Wegen $\Gamma(\nu+k+1) = \Gamma(\nu+1)(\nu+1)(\nu+2)\cdots(\nu+k)$ verschluckt die Gammafunktion genau das Produkt aus dem Nenner von~\eqref{bessel:loesung:eqn:a2k}.
Damit erhält die Lösung die Gestalt
\begin{definition}
\label{bessel:loesung:def:jnu}
Die \emph{Bessel-Funktion erster Art} der Ordnung $\nu$
\index{Bessel-Funktion erster Art}%
ist
\begin{equation}
J_\nu(x)
=
\sum_{k=0}^\infty
\frac{(-1)^k}{k!\,\Gamma(\nu+k+1)}
\Bigl(\frac{x}{2}\Bigr)^{2k+\nu}.
\label{bessel:loesung:eqn:jnu}
\end{equation}
\end{definition}
 
Die Reihe~\eqref{bessel:loesung:eqn:jnu} konvergiert für alle $x$.
Für ganzzahliges $\nu=n$ ist $\Gamma(n+k+1)=(n+k)!$, und man erhält
\[
J_n(x)
=
\sum_{k=0}^\infty
\frac{(-1)^k}{k!\,(n+k)!}
\Bigl(\frac{x}{2}\Bigr)^{2k+n}.
\]
Besonders anschaulich ist der Fall $n=0$,
\[
J_0(x)
=
1
-
\frac{x^2}{4}
+
\frac{x^4}{64}
-
\frac{x^6}{2304}
+
\dots,
\]
der an die Reihe des Kosinus erinnert, aber schneller abfällt.
% [TODO: Abbildung mit J_0, J_1, J_2 auf dem Intervall [0,20],
%        pgfplots, Nullstellen markieren]
Für kleine Argumente dominiert der erste Term,
\begin{equation}
J_\nu(x)
\approx
\frac{1}{\Gamma(\nu+1)}
\Bigl(\frac{x}{2}\Bigr)^{\nu}
\qquad (x\to 0),
\label{bessel:loesung:eqn:kleinx}
\end{equation}
insbesondere ist $J_0(0)=1$ und $J_\nu(0)=0$ für $\nu>0$.
 
%
% Die zweite Lösung
%
\subsection{Die zweite Lösung
\label{bessel:loesung:subsection:ynu}}
Eine Differentialgleichung zweiter Ordnung hat einen zweidimensionalen Lösungsraum, und bisher ist erst eine Lösung gefunden.
Die Indexgleichung~\eqref{bessel:loesung:eqn:index} hat aber eine zweite Wurzel $\varrho_2=-\nu$, und die ganze Rechnung lässt sich mit diesem Exponenten wiederholen.
Man erhält
\begin{equation}
J_{-\nu}(x)
=
\sum_{k=0}^\infty
\frac{(-1)^k}{k!\,\Gamma(-\nu+k+1)}
\Bigl(\frac{x}{2}\Bigr)^{2k-\nu}.
\label{bessel:loesung:eqn:jminusnu}
\end{equation}
 
Ob damit tatsächlich eine zweite, unabhängige Lösung vorliegt, hängt vom Wert von $\nu$ ab.
Ist $\nu\notin\mathbb{Z}$, so verhält sich $J_\nu$ bei $x=0$ wie $x^{\nu}$ und $J_{-\nu}$ wie $x^{-\nu}$.
Die eine Funktion bleibt beschränkt, die andere nicht, sie können also nicht proportional sein.
In diesem Fall ist
\[
y(x)
=
C_1 J_\nu(x) + C_2 J_{-\nu}(x)
\]
die allgemeine Lösung.
Für ganzzahliges $\nu=n$ bricht dieses Argument zusammen.
Die Gammafunktion hat an den nichtpositiven ganzen Zahlen Pole, ihr Kehrwert verschwindet dort.
In der Reihe~\eqref{bessel:loesung:eqn:jminusnu} fallen deshalb die ersten $n$ Glieder weg, und eine Indexverschiebung ergibt
\begin{equation}
J_{-n}(x) = (-1)^n J_n(x).
\label{bessel:loesung:eqn:jminusn}
\end{equation}
% [TODO: Indexverschiebung k -> k+n explizit vorrechnen]
Die zweite Wurzel der Indexgleichung liefert also nichts Neues.
Das ist kein Zufall, sondern der im Fuchs-Abschnitt beschriebene Ausnahmefall.
Die beiden Wurzeln $\varrho_1=n$ und $\varrho_2=-n$ unterscheiden sich um die ganze Zahl $2n$, und dann kann die Rekursion für den kleineren Exponenten zusammenbrechen.
Man definiert deshalb
\begin{definition}
\label{bessel:loesung:def:ynu}
Die \emph{Bessel-Funktion zweiter Art} der Ordnung $\nu$
\index{Bessel-Funktion zweiter Art}%
oder \emph{Neumann-Funktion}
\index{Neumann-Funktion}%
ist
\begin{equation}
Y_\nu(x)
=
\frac{J_\nu(x)\cos(\nu\pi) - J_{-\nu}(x)}{\sin(\nu\pi)}
\label{bessel:loesung:eqn:ynu}
\end{equation}
für $\nu\notin\mathbb{Z}$ und
\begin{equation}
Y_n(x)
=
\lim_{\nu\to n} Y_\nu(x)
\label{bessel:loesung:eqn:yn}
\end{equation}
für $n\in\mathbb{Z}$.
\end{definition}
 
Für nicht ganzzahliges $\nu$ ist $Y_\nu$ nur eine Linearkombination von $J_\nu$ und $J_{-\nu}$ und damit als Lösung nichts Neues.
Der Sinn der Definition liegt darin, dass der Grenzwert \eqref{bessel:loesung:eqn:yn} existiert und dort eine von $J_n$ unabhängige Lösung liefert, wo $J_{-n}$ versagt.
Zähler und Nenner von~\eqref{bessel:loesung:eqn:ynu} streben für $\nu\to n$ beide gegen null, und die Regel von de l'Hôpital liefert einen endlichen Grenzwert.
% [TODO: den Grenzübergang wenigstens für Y_0 durchrechnen oder
%        auf Watson verweisen]
 
Dieser Grenzwert enthält einen logarithmischen Term.
Für $\nu=0$ etwa ist
\begin{equation}
Y_0(x)
=
\frac{2}{\pi}
\Bigl(
\ln\frac{x}{2} + \gamma
\Bigr)
J_0(x)
+
\frac{2}{\pi}
\sum_{k=1}^\infty
\frac{(-1)^{k+1}}{(k!)^2}
\Bigl(\frac{x}{2}\Bigr)^{2k}
\sum_{m=1}^k \frac{1}{m},
\label{bessel:loesung:eqn:y0}
\end{equation}
wobei $\gamma$ die Euler-Mascheroni-Konstante bezeichnet.
Der Logarithmus lässt sich durch keine verallgemeinerte Potenzreihe darstellen, und deshalb konnte der Ansatz \eqref{bessel:loesung:eqn:ansatz} diese Lösung nicht finden.
Insbesondere gilt $Y_\nu(x)\to-\infty$ für $x\to 0$.
Damit ist die allgemeine Lösung der Bessel-Gleichung für jede Ordnung
\begin{equation}
y(x)
=
C_1 J_\nu(x) + C_2 Y_\nu(x).
\label{bessel:loesung:eqn:allgemein}
\end{equation}
In Anwendungen auf ein Gebiet, das den Ursprung enthält, etwa bei der kreisförmigen Membran, muss $C_2=0$ gesetzt werden, weil die Lösung dort beschränkt bleiben muss.
Bei Ringgebieten oder im Aussenraum wird $Y_\nu$ dagegen gebraucht.